\documentclass[12pt]{amsart}
\input{../../../../../Headers/TestPackageHeader}

%\fancyfoot[CE]{\begin{tabular}{|c|} \hline $\quad\qquad$ \\ $\quad\qquad$ \\ $\quad\qquad$ \\ \hline \end{tabular}}
%\fancyfoot[RO]{\begin{tabular}{|c|} \hline $\quad\qquad$ \\ $\quad\qquad$ \\ $\quad\qquad$ \\ \hline \end{tabular}}
\renewcommand{\headrulewidth}{0pt} % remove lines as well

\numberwithin{equation}{section}

\DeclareGraphicsRule{.tif}{png}{.png}{`convert #1 `dirname #1`/`basename #1 .tif`.png}

\textheight = 10.9 in
\topmargin = -1.1in
\textwidth = 7.7in
\oddsidemargin = -0.7in
\evensidemargin = -0.7in

\begin{document}

\input{../../../../../Headers/TestDefinitionHeader.tex}

\vspace{1.0cm}

\begin{center}

$\vspace{1.2cm}$

\LARGE {MA 301} \\

\Large{Final} \\

\large{Fall - 2011} \\

\vspace{0.5cm}

\setlength{\extrarowheight}{0.1cm}
\begin{tabular}{|c|c|c|}
\hline
Problem & points & out of  \\
\hline
\hline
1 & & 10 \\
\hline
2 & & 10 \\
\hline
3 & & 10 \\
\hline
4 & & 10 \\
\hline
5 & & 10 \\
\hline
6 & & 10 \\
\hline
7 & & 10 \\
\hline
\hline
Score & & 70\\
\hline
\end{tabular}
\vspace{0.5cm}

%--------TABLE---------
	\setlength{\extrarowheight}{0.2cm}
	\begin{tabular}{rcc}
		Name: & \underline{$\qquad \qquad \qquad \qquad \quad$}& \\
		&&$\hspace{1.4cm}$\\
		HR: & \underline{$\qquad \qquad \qquad \qquad \quad$}& \\
		&&\\
		Signature: & \underline{$\qquad \qquad \qquad \qquad \quad$}& \\
	\end{tabular}
%--------------------

\vspace{0.5cm}

	\setlength{\extrarowheight}{0.2cm}
	\begin{tabular}{| p{12cm} |}
		\hline
		\multicolumn{1}{|c|}{{\bf READ BEFORE YOU START}} \\
		\hline
		 * Show all work.\\
		 * Change your calculator mode to RADIANS.\\
		 * Sloppy solutions may be penalized.\\
		 * Use calculators ONLY for computations.\\
		 * You must write up every problem in a clear and \\
		 $\,\,$ organized way.\\
		 * If you write solutions on a page different than the \\
		 $\,\,$ question, then clearly state where it is located.\\
		\hline
	\end{tabular}

\end{center}

\vspace{1.0cm}

\thispagestyle{empty}
\newpage
\thispagestyle{empty}

{\bf Problem 1:} Given $\text{\bf F} = [2,2y+z,y+1]$

(a) Show that $\text{\bf F}$ is conservative in two different ways

	\begin{tabular}{p{10cm}|p{10cm}}
		\multicolumn{1}{c|}{method 1}&\multicolumn{1}{c}{method 2}\\
		&\\&\\&\\&\\&\\&\\&\\&\\&\\&\\&\\&\\&\\&\\&\\&\\&\\&\\&\\&\\&\\&\\&\\&\\&\\&\\&\\&\\&\\&\\&\\&\\&\\&\\&\\&\\&\\
	\end{tabular}

\vspace{0.4cm}

(b) Evaluate $\displaystyle{\int \limits_{C} \text{\bf F} \cdot  d\text{\bf r}}$ when $C$ is a straight line from $(-1,0,2)$ to $(0,2,-6)$

\thispagestyle{empty}
\newpage
\thispagestyle{empty}

{\bf Problem 2:} Let $\text{\bf F} = [y,x^2\sin y]$ and $C$ be the boundary of the region bounded by $\D{y=\frac{\pi}{2}x}$, \, $\D{y=\frac{\pi}{2}}$ \, and $0\le x \le 1$ \, in the $x,y$-plane.

Sketch $C$ and compute \,\, $\D{\oint \limits_{C} \text{\bf F}\cdot d\text{\bf r}}$ \,\, using green's theorem.\\

\vspace{19.2cm}

\thispagestyle{empty}
\newpage
\thispagestyle{empty}

{\bf Problem 3:} Let $\text{\bf F} = [e^x,e^y,e^z]$ and $S$ be the surface of the box $|x|\le 1$, $|y|\le 2$ , $0\le z \le 1$.

\vspace{0.2cm}

Compute \,$\D{\iint \limits_{S} \text{\bf F}\cdot \text{\bf n} \, dA}$

\vspace{0.2cm}

\thispagestyle{empty}
\newpage
\thispagestyle{empty}

{\bf Problem 4:} Answer the following questions

\vspace{0.1cm}

(a) If the integral \,\, $\D{\oint \limits_{C} \text{\bf F}\cdot d\text{\bf r}}$ \,\, represents the work done by a object then explain the physical role of the vector-valued function {\bf F} and the curve $C$.

\vspace{5.0cm}

(b) \quad Suppose $f(x) = f(x+6)$ for all $x$, then we say $f$ is periodic with period 3.  \, TRUE \, or \, FALSE

\vspace{2.0cm}

(c) \quad $\sqrt{-5}\cdot\sqrt{-5} = 5$  \quad TRUE \, or \, FALSE

\vspace{2.0cm}

(d) Let $\text{\bf F}$ vector field such that curl\,$\text{\bf F} =$ {\bf 0}.  Match each expression in the left column with the most appropriate term in the right column.

\begin{center}
%\setlength{\extrarowheight}{0.8cm}
\begin{tabular}{cp{5cm}c}
	$\D{\int \limits_{C} \text{\bf F}\cdot d\text{\bf r}}$ & & conservative\\
	&&\\
	$\D{\text{\bf F}\cdot d\text{\bf r}}$ & & zero\\
	&&\\
	$\D{\oint \limits_{C} \text{\bf F}\cdot d\text{\bf r}}$ & & path independent\\
	&&\\
	$\D{\text{\bf F}}$ & & exact\\
\end{tabular}
\end{center}

\vspace{2.0cm}

(e) Match each the Theorem in the left column with the most appropriate statement in the right column.

\begin{center}
%\setlength{\extrarowheight}{0.8cm}
\begin{tabular}{cp{5cm}c}
	Divergence Theorem & & The surface of a sphere\\
	&&\\
	Stokes' Theorem & & Top-half of the surface of a sphere\\
	&&\\
	Green's Theorem & & The surface of a circle\\
\end{tabular}
\end{center}

\thispagestyle{empty}
\newpage
\thispagestyle{empty}

{\bf Problem 5:} Let \quad
$
f(x) = 
\left\{
\begin{array}{lrc}
	0 & &-2  < x <  -1\\
	1+x & &-1  < x <   0\\
	1-x & &0  < x <   1\\
	0 & & 1  < x <   2\\
\end{array}
\right.
$  \quad and \quad $f(x) = f(x+4)$

\vspace{0.1cm}

(a) Graph $f(x)$ on the interval $[-3,5]$

\vspace{7.2cm}

(b) Find the Fourier Series of $f(x)$

\thispagestyle{empty}
\newpage
\thispagestyle{empty}

{\bf Problem 6:} Let \quad 
$\D{f(x) = 
\left\{
\begin{array}{rrr}
	-x^2 & &-1  < x <   0\\
	x^2 & &0  < x <   1\\
\end{array}
\right.
\quad \text{ and } \quad f(x) = f(x+2)}$

\vspace{0.1cm}

(a) Graph $f(x)$ on the interval $[-3,3]$

\vspace{8.2cm}

(b) Can we say that any of Fourier coefficients $a_0$, $a_n$, or $b_n$ are zero? Why or why not?

\vspace{2.2cm}

(c) {\bf Set-up}, but {\bf DO NOT} compute the integral(s) for all the non-zero Fourier coefficients.

\thispagestyle{empty}
\newpage
\thispagestyle{empty}

{\bf Problem 7:} Let \,\,$z_1 = -2-2i$,\,\, $z_2 = i$. \quad Answer the following:

\vspace{0.2cm}

(a) Compute the operations and write the result in standard form 

\vspace{0.2cm}

\hspace{0.6cm} (i) $z_1z_2 = $

\vspace{2.8cm}

\hspace{0.6cm} (ii) $z_1-z_1z_2 = $

\vspace{3.8cm}

\hspace{0.6cm} (iii) $|z_1z_2| = $

\vspace{3.8cm}

\hspace{0.6cm} (iv) $(z_2)^{76}$

\vspace{3.8cm}

(d) Solve for $z$: \quad $\D{z_2z^{1/3} = z_1}$

 
\thispagestyle{empty}
\newpage
\thispagestyle{empty}

{\bf Problem 1:} Line Integral - Directly

$\text{\bf F} = [2,2y+z,y+1]$  

\vspace{0.2cm}

%{\bf Problem 2:} Line Integral - Potential
%
%%$\D {\bf F} = \left[2xy,x^2-\frac{z}{y},-2yz\right]$
%
%\includegraphics[scale = 0.3]{../../../Figures/Test1_Fig2.png}
%
%\vspace{0.2cm}
%
%{\bf Problem 3:} Line Integral - Green's Thm
%
%Let ${\bf F} = [x^2 + 4y, x + y^2]$ and $C$ be the rectangle in the $xy-$plane with $0 \le x \le 4$ and $-1 \le y \le 2$

\vspace{0.2cm}

Compute \quad $\D \OWORK $

Let ${\bf F} = [-y,x]$ and $C$ be a simple closed curve in the $xy-$plane. If $R$ is the region inside $C$, show that 
$$
\text{Area}(R) = \frac{1}{2}\oint \limits_{C} \text{\bf F} \cdot d\text{\bf r}
$$

Use part (a) to determine the area inside the ellipse $C$: $4x^2 + y^2 = 4$

\vspace{0.2cm}

{\bf Problem 4:} Flux Integral - Directly

\vspace{0.2cm}

{\bf Problem 5:} Flux Integral - Directly with unit normal

\vspace{0.2cm}

{\bf Problem 6:} Flux Integral - Divergence Thm

\vspace{0.2cm}

{\bf Problem 7:} Conceptual

\EX If $\text{\bf F} = [y^2,-x^2]$ and $C: y = 4x^2$ \,\, from \, $(0,0)$ \, to \, $(1,4)$, then compute $\D{ \int \limits_{C} \text{\bf F} \cdot d\text{\bf r}}$

\vspace{0.2cm}

\EX	Determine $\D{ \int \limits_{C} \text{\bf F} \cdot d\text{\bf r}}$ \, where \, $\text{\bf F} = [x,y-z^2,z-x]$ and $C: \text{\bf r}(t) = [2\cos t, t, 2\sin t]$ from $(2,0,0)$ to $(2,2\pi,0)$

\vspace{0.2cm}

$\text{\bf F} = [x,1]$  with \, $\text{\bf r}(t) = [\cos t, \sin t]$ \, \& \, $\D{0 \le t \le \frac{\pi}{2}}$
	
\vspace{0.2cm}
	
What is the work done to move the object along this path when exposed to \text{\bf F}?
	
\vspace{0.2cm}

\EX Let $\D\text{\bf F} = [xy,2y^2]$ represent a force vector field in the x-y plane.  If a particle moves along a quarter circle from $(2,0)$ \, to \, $(0,2)$, then find the work done on that particle.

\vspace{0.2cm}

\EX Suppose a wire has density $\delta(x,y,z) = 3\sqrt{5+x}$.  As a curve, the is parameterized by \RT = t\VI + 2t\VJ + $\frac{2}{3}t^{3/2}$\VK \,where\, $0 \le t \le 2$. Find the mass of the wire.

\vspace{0.2cm}

\EX Let $\D{{\bf F} = \nabla(x^2\sin y)}$ and let\, $R$ be the region bounded by: \,\, $x=1$, \, $y=0$, \, and \, $y=\sqrt{x}$

\vspace{0.2cm}

Compute \quad $\D{\OWORK}$

\vspace{0.2cm}

\EX Compute the flux of water through the parabolic cylinder $S:\, y=x^2$ where $0 \le x \le 2$ and $0 \le z \le 3$ if the velocity flow of water is given by ${\bf v} = [3z^2, 6, 6xz]$ (units: {\em m/s}).  Note, the density of water is $\rho = 1 \, kg/m^3$.

\vspace{0.2cm}

\EX Compute the surface integral of ${\bf F} = [e^y, -e^z, e^x]$ over the surface $S:\, x^2 + y^2 = 25$ \, for $x \ge 0$, $y \ge 0$, and $0 \le z \le 2$

\vspace{0.2cm}

\EX Compute the surface integral \qquad $\D{\iint \limits_{S} (1 + 9xz)^{3/2}\, dA}$

\vspace{0.1cm}

where $S$ is given by the parameterization: ${\bf r}(u,v) = [u,v,u^3]$,\,\, $0 \le u \le 1$,\,\, $-2 \le v \le 2$

\vspace{0.2cm}

Compute the flux of ${\bf F} = [x^2,0,z^2]$ \, through the rectangular box: $|x| \le 1$, \,\, $|y| \le 3$, \,\, and $0 \le z \le 2$.

\vspace{0.2cm}

\EX Evaluate \quad $\D{\iint \limits_{S} {\bf F} \cdot {\bf n} \, dA}$	\quad if $F = [x^3-y^3,y^3-z^3,z^3-x^3]$ and $S$ is given by: $x^2 + y^2 + z^2 \le 25$, \, $z \ge 0$

\vspace{0.2cm}

The closed circular cylinder shown below has a height of 10 and a base radius of 6. Let $A$ be the top, $B$ the bottom, and $S$ the side surface of the cylinder. Also, let ${\bf F} = (7x+1){\bf i} - 2y{\bf j} + z{\bf k}$. Find the outward flux through side $S$.  ({\em Note}:  $Vol(cylinder) = \pi r^2 h$ and $Area(circle) = \pi r^2$)	

\vspace{0.2cm}

{\bf Problem 2:} Let $\text{\bf F} = [x+y,2xy]$ and $C$ be the boundary of the region bounded by $y=x^2$ and $y=x+2$ in the $x,y$-plane. \\
(a) Compute \,\, $\D{\oint \limits_{C} \text{\bf F}\cdot d\text{\bf r}}$ \,\, using green's theorem.\\
(b) \UL{Setup} the integral you would use to compute \,\, $\D{\oint \limits_{C} \text{\bf F}\cdot d\text{\bf r}}$ \,\, directly \,\,\,({\em don't leave answer as a dot product})

\vspace{0.2cm}

{\bf Problem 3:} Let $\text{\bf F} = [x,y,z]$ and $S$ be the sphere $x^2+y^2+z^2 = 25$.  \,\,Compute \,$\D{\iint \limits_{S} \text{\bf F}\cdot \text{\bf n} \, dA}$

\vspace{0.2cm}

{\bf Problem 4:} Compute the flux of water (water density $= 1 \, ton/m^3$) through the parabolic cylinder $S:\, y=x^2$ where $0 \le x \le 2$ and $0 \le z \le 1$ if the velocity flow of water is given by $\text{\bf v} = [3z^2, 6, 6xz]$ (units: {\em m/s}).

\vspace{0.2cm}

\newpage 

{\bf Problem 5:} Answer the following conceptual questions

\vspace{0.1cm}

(a) \quad Suppose $\text{\text{\bf F}} = [1,1,1]$.  Is it true that \,\,$\D{\iint \limits_{S} \text{\bf F}\cdot \text{\bf n} \, dA = 0}$ \,\,for any closed surface $S$ which is a boundary of a solid?

\vspace{0.2cm}

The line integral of any vector field {\bf F} around any closed curve $C$ is zero

$\WORK = f(Q) - f(P)$ whenever $P$ and $Q$ are the endpoints of the curve $C$.

If $\WORK = 0$ for one particular closed path, the ${\bf F}$ is path-independent.

If ${\bf F}$ is path-independent, then there is a potential function for ${\bf F}$ 

If $\WORK \ne 0$ for some closed path $C$, then ${\bf F}$ is path-dependent.

If ${\bf F} = \nabla f$ then ${\bf F}$ is conservative.

If ${\bf F}$ is path-independent and $C$ is any closed curve, then $\WORK = 0$

If $S_1$ is a rectangle with area = 1 and $S_2$ is a rectangle with area = 2, then
$2\iint \limits_{S_1} {\bf F} \cdot d{\bf A} = \iint \limits_{S_2} {\bf F} \cdot d{\bf A}$

Let $S$ be a sphere of radius 1 centered at the origin and oriented outward.  If ${\bf F} = [x,y,z]$ then $\FLUX$ is positive

Let $S$ be a sphere of radius 1 centered at the origin and oriented outward.  If $\FLUX = 0$, then ${\bf F} = {\bf 0}$

Let $S$ be the surface of a cube centered at the origin with side length 2 and oriented outward.  If ${\bf F} = 2\VI - \VK$ then the flux of ${\bf F}$ through $S$ is 27.

If ${\bf F}$ and ${\bf G}$ are vector-fields satisfying div${\bf F}$ = div${\bf G}$, then ${\bf F} = {\bf G}$

If ${\bf F}$ is a vector-field such that div{\bf F} = 1 and $S$ is a closed surface oriented outward, then $\FLUX$ is equal to the volume enclosed by $S$.

If $\FLUX = 12$ and $S$ is a flat disk of area $4\pi$, then div${\bf F} = \frac{3}{\pi}$ 
\thispagestyle{empty}
\newpage
\thispagestyle{empty}

{\bf Problem 5:} Answer the following true/false questions

\vspace{0.1cm}

	\setlength{\extrarowheight}{0.2cm}
	\begin{tabular}{| c | p{17.2cm} |}
		\hline
		T\quad/\quad F & The line integral of any vector field {\bf F} around any closed curve $C$ is zero\\
		\hline
		T\quad/\quad F & If $\WORK = 0$ for one particular closed path, the ${\bf F}$ is path-independent.\\
		\hline
		T\quad/\quad F & If ${\bf F}$ is path-independent, then there is a potential function for ${\bf F}$\\
		\hline
		T\quad/\quad F & If $\WORK \ne 0$ for some closed path $C$, then ${\bf F}$ is path-dependent.\\
		\hline
		T\quad/\quad F & If ${\bf F} = \nabla f$ then ${\bf F}$ is conservative.\\
		\hline
		T\quad/\quad F & If ${\bf F}$ is path-independent and $C$ is any closed curve, then $\WORK = 0$\\
		\hline
		T\quad/\quad F & If $S_1$ is a rectangle with area = 1 and $S_2$ is a rectangle with area = 2, then
$2\iint \limits_{S_1} {\bf F} \cdot d{\bf A} = \iint \limits_{S_2} {\bf F} \cdot d{\bf A}$\\
		\hline
		T\quad/\quad F & Let $S$ be a sphere of radius 1 centered at the origin and oriented outward.  If ${\bf F} = [x,y,z]$ then $\FLUX$ is positive\\
		\hline
		T\quad/\quad F & Let $S$ be a sphere of radius 1 centered at the origin and oriented outward.  If $\FLUX = 0$, then ${\bf F} = {\bf 0}$\\
		\hline
		T\quad/\quad F & Let $S$ be the surface of a cube centered at the origin with side length 2 and oriented outward.  If ${\bf F} = 2\VI - \VK$ then the flux of ${\bf F}$ through $S$ is 27.\\
		\hline
		T\quad/\quad F & If ${\bf F}$ and ${\bf G}$ are vector-fields satisfying div${\bf F}$ = div${\bf G}$, then ${\bf F} = {\bf G}$\\
		\hline
		T\quad/\quad F & If ${\bf F}$ is a vector-field such that div{\bf F} = 1 and $S$ is a closed surface oriented outward, then $\FLUX$ is equal to the volume enclosed by $S$.\\
		\hline
		T\quad/\quad F & If $\FLUX = 12$ and $S$ is a flat disk of area $4\pi$, then div${\bf F} = \frac{3}{\pi}$\\
		\hline
	\end{tabular}

\end{document}